% Chapter 3: Rounding and the choice of step (fundamentals author). PLAN.md section 1.4.
% Uses F3 (fig:v-curve), F4 (fig:step-scale), F4b (fig:step-scale-small), F5 (fig:float-double), F17 (fig:noise),
% S3 (sb:ulp), S4 (sb:actual-distance), S5 (sb:step-constant), E3 (code/03-rounding-step.cpp),
% E4 (code/03-rounding-scale.cpp), E17 (code/03-rounding-noise.cpp).
% Language pass for students (PLAN.md section 12.9): the wording only; content, evidence and order as at 3192e8a.
\chapter{Rounding and the choice of step}\label{ch:rounding}

\Cref{ch:taylor} predicts that a smaller step gives a smaller error. Halve $h$, and the error of the central
difference falls by a factor of four. A computer does not follow this prediction for long. To see it, pass the step
yourself, as the third argument of \code{derivative}, and compare a moderate step with a tiny one:
\snippet{03-rounding-step}{explicit-step}
The tiny step of \SI{e-8}{\kelvin} is not better than the step of \SI{e-2}{\kelvin}. It is worse, by orders of
magnitude. This chapter explains why, and then finds the best step. It also shows how the library builds its default
step, when that step is wrong, and what to do then.

\section{Floating-point numbers in brief}\label{sec:rounding-fp}

To explain the result, we need a few facts about how a computer stores numbers. A \code{double} stores a number with
\num{\valDigitsDouble} binary digits. That is about 16 significant decimal digits. So a \code{double} cannot hold
every real number: the numbers it can hold are spaced apart. Just above 1, the next \code{double} is $1 + \epsmach$,
where
\begin{equation*}
  \epsmach = 2^{-52} \approx \num{\valEpsDouble}
\end{equation*}
is the \emph{machine epsilon} of the type.\mnote{\code{std::numeric_limits<T>::epsilon()} in C++.} The spacing grows
in proportion to the numbers themselves. Near \num{1000}, for example, the doubles are \num{\valUlpAtOneThousand}
apart.

Every arithmetic operation rounds its exact result to the nearest number that the type can represent. This changes
the result by at most half a spacing, which is a relative error of at most
$\epsmach / 2$.\cite{higham2002,goldberg1991} Now suppose that a value of $f$ is computed by a few dozen operations. Each of them
can add a rounding error. So the value is wrong by a few units of $\epsmach$, relative to $f$: sometimes less,
sometimes more. We model all of this with one number. Each value of $f$ carries an error of about
$\epsmach \abs{f}$.

That is all you need to know about floating point to follow this chapter. \Cref{ch:floating-point} has the details.
\Cref{sb:ulp}, in \cref{sec:rounding-default}, explains what the spacing of the numbers near $x$ means for a step.

\section{The two errors and the V}\label{sec:rounding-v}

We can now look at what goes wrong. \Cref{fig:v-curve} shows the error of the central difference for the heat
capacity of nitrogen, over fifteen decades of step (a decade is a factor of ten). Read the figure from the right.
The error first falls along a line of slope 2, as \cref{ch:taylor} predicts, down to a step of about
\SI{0.01}{\kelvin}. From there to the left it \emph{rises} again, along a line of slope $-1$. At the smallest steps
the result has no correct digit at all. The curve has the shape of a V.

\begin{figure}
\sidecaption{The V. The relative error of \code{Central<2>} against the step $h$, for the heat capacity of nitrogen at
$T_0 = \SI{1000}{\kelvin}$, over the whole range of steps. The dashed line is the model
$E(h) / \cp$ of \cref{eq:error-model}. To the right of the bottom, truncation dominates and the error falls as $h^2$.
To the left of it, rounding dominates and the error grows as $1/h$. The default step sits near the
bottom.\label{fig:v-curve}}
\centering
% F3 fig:v-curve (Chapter 3, text width): the V of the central difference (Central<2>) on the nitrogen example, H(T) at
% T0 = 1000 K: the measured relative error over the whole range of steps and the model
% E(h) = (h^2 |H'''| / 6 + eps |H| / h) / cp, dashed (DS1, data/step-error-cp.csv); the default step marked from
% \valCpDefaultStep and \valCpDefaultRelError (DS1b). The model exceeds 1 at the smallest steps, cut at the top.
\begin{tikzpicture}
\begin{axis}[eqtufte, width=95mm, height=45mm, xmode=log, ymode=log,
    xlabel={step $h$ (K)}, ylabel={relative error},
    xmin=1e-13, xmax=1e2, ymin=1e-12, ymax=1,
    xtickten={-12,-10,-8,-6,-4,-2,0,2}, ytickten={-12,-9,-6,-3,0},
    restrict y to domain=-30:0.01, unbounded coords=jump]
  \addplot[muted, densely dashed, line width=0.6pt] table[col sep=comma, x=h, y=model_central2,
    comment chars={\#}] {data/step-error-cp.csv};
  \addplot[accent, mark=*] table[col sep=comma, x=h, y=central2, comment chars={\#}] {data/step-error-cp.csv};
  % the default step: a thin muted line over the height of the axis, labelled at the top
  \draw[muted, thin] ({axis cs:\valCpDefaultStep,1}|-{rel axis cs:0,0})
    -- ({axis cs:\valCpDefaultStep,1}|-{rel axis cs:0,1})
    node[above, font=\sffamily\scriptsize, inner sep=1.5pt] {default step};
  \fill[accent] (axis cs:\valCpDefaultStep,\valCpDefaultRelError) circle (1.4pt);
  % the two branches and the model
  \node[font=\sffamily\scriptsize, black!70, anchor=south west, align=left] at (axis cs:3e-13, 1e-10)
    {rounding:\\ error ${}\propto 1/h$};
  \node[font=\sffamily\scriptsize, black!70, anchor=north west, align=left] at (axis cs:0.3, 1e-9)
    {truncation:\\ error ${}\propto h^2$};
  \node[font=\sffamily\scriptsize, muted, anchor=south west] at (axis cs:1e-9, 1e-3) {model $E(h)$, dashed};
\end{axis}
\end{tikzpicture}

\end{figure}

The right branch is the truncation error of \cref{eq:central-error}. The left branch comes from rounding. To see how,
look at what the central difference computes. It subtracts $f(x - h)$ from $f(x + h)$. As $h$ shrinks, these two
numbers agree in more and more of their leading digits. The subtraction itself is exact. But it removes the digits
that the two values share, and leaves only the digits in which they differ. Those digits include the rounding errors
of both values. The rounding errors were tiny compared with $f$, but they are not tiny compared with the difference.
This loss of digits is called \emph{cancellation}.

You can watch this happen on a pocket calculator that keeps six significant digits. Four lines show the whole V. We
differentiate $\mathrm{e}^x$ at $x = 1$, where $f'(1) = \mathrm{e} = \num{\valHandCalcExact}\ldots$, and round each
value of $f$ to six digits:
\begin{center}\small
\begin{tabular}{@{}lllll@{}}
\toprule
$h$ & $f(1 + h)$ & $f(1 - h)$ & estimate & error\\
\midrule
\tablebody{hand-calculator}
\bottomrule
\end{tabular}
\end{center}
From $h = 0.1$ to $0.01$ the step shrinks by a factor of 10. Since the truncation error falls as $h^2$, truncation
alone would make the error 100 times smaller. Instead, the error falls by a factor of only
\num{\valHandCalcErrorRatio}. The reason is that at $h = 0.01$ the truncation error is only about
\num{\valHandCalcTruncationHundredth}. The rest of the error in that row is already rounding. Below that step, the
error rises. At $h = 10^{-4}$ the two values share their first four digits. Their difference,
\num{\valHandCalcLastDiff}, has only two digits left, and the last of them is rounding. A \code{double} keeps about 16
digits instead of six. So the same thing happens at smaller steps, but it does happen.

We can write both errors down. Suppose that each value of $f$ is off by up to about $\epsmach \abs{f}$. In the worst
case the two errors have opposite signs and add up. So the difference $f(x + h) - f(x - h)$ is off by up to about
$2 \epsmach \abs{f}$. After the division by $2h$, the estimate of $f'$ is off by up to about
$2 \epsmach \abs{f} / (2h) = \epsmach \abs{f} / h$. Adding the truncation error of \cref{eq:central-error} gives the
\emph{error model} of the central difference,
\begin{equation}
  E(h) = \etrunc(h) + \eround(h) = \frac{h^2}{6} \abs{f'''(x)} + \frac{\epsmach \abs{f(x)}}{h}.
  \label{eq:error-model}
\end{equation}
As $h$ shrinks, the first term, the truncation error, falls, and the second, the rounding error, grows. Their sum is
the V. The dashed line in \cref{fig:v-curve} is this model for the nitrogen example. The measured curve follows it,
with one difference. On the rounding branch, the measured errors scatter below the model. This is because the
rounding errors of the two values have random signs, and they partly cancel. So the model gives their typical
largest size, not their actual size at a particular step.

Two more remarks on the model. First, it gives an \emph{absolute} error, in the unit of $f'$. The figure divides it by
$\cp$, to show a relative error. Second, it leaves out a second rounding: that of the points $x \pm h$ themselves. For
the central difference, the library removes this rounding by dividing by the distance actually evaluated
(\cref{sb:actual-distance}).

\begin{sidebar}[label=sb:actual-distance]{Dividing by the distance actually evaluated}
The points $x + h$ and $x - h$ are computed, so they are rounded too. The function is evaluated at
$x_+ = \fl(x + h)$ and $x_- = \fl(x - h)$, the rounded values of the two sums. Each lies within half a spacing of the
intended point. So their distance can differ from $2h$ by up to about one spacing, $\epsmach \abs{x}$. Suppose that we
divided by the nominal $2h$ anyway. That would add a relative error of about $\epsmach \abs{x} / (2h)$. At the
default step $h = \epsmach^{1/3} \abs{x}$, this is $\epsmach^{2/3} / 2$: as large as the accuracy we aim for.

This is why the two-node rules (\code{Central<2>}, \code{Forward<1>} and \code{Backward<1>} for $f'$) divide by
$x_+ - x_-$ as evaluated. The quotient is then the exact slope of the secant through the two points where $f$ was
actually evaluated. Only the rounding of the values of $f$ remains. The subtraction $x_+ - x_-$ is itself exact
when $x_+ \le 2 x_-$ (for positive $x$). This follows from Sterbenz's lemma (Theorem 4.3.1):\cite{sterbenz1974} two
floating-point numbers of the same sign, within a factor of two of each other, have an exactly representable
difference (Higham's Theorem 2.5).\cite{higham2002} Since $x + h \le 2(x - h)$ means $3h \le x$, the condition holds
roughly for $h \le \abs{x} / 3$. For larger steps the subtraction rounds once. That is a relative error of at most
$\epsmach / 2$, which is harmless.

Rules with more nodes divide by a fixed multiple of the nominal step instead. For them, the actual distance would not
help. Take \code{Central<4>}. Dividing by the distance between its outer nodes would quadruple the point rounding of
those nodes (its coefficient rises from $1/12$ to $4/12$), and would leave that of the inner nodes untouched. So these
rules keep a rounding term from the points, of about
$(\ampl / 2)\, \epsmach \abs{x f'} / h$, where $\ampl$ is the sum of the absolute weights (\cref{ch:stencils}).
\end{sidebar}

\section{The best step and the best accuracy}\label{sec:rounding-best}

We now know where the two errors come from. Next we find the step that makes their sum smallest. The bottom of the V
is where the model is smallest, and there its derivative with respect to $h$ is zero. Differentiating
\cref{eq:error-model} with respect to $h$ and setting the result to zero gives
$h \abs{f'''} / 3 - \epsmach \abs{f} / h^2 = 0$, so $h^3 = 3 \epsmach \abs{f} / \abs{f'''}$. This is the \emph{best
step}, with the model error that goes with it:
\begin{equation}
  \hbest = \left( \frac{3 \epsmach \abs{f}}{\abs{f'''}} \right)^{1/3},
  \qquad
  E(\hbest) = \frac{3}{2} \frac{\epsmach \abs{f}}{\hbest}.
  \label{eq:best-step}
\end{equation}
The second formula follows from putting $\abs{f'''} = 3 \epsmach \abs{f} / (\hbest)^3$ into the model: at the best step,
the truncation term is half the rounding term.

Both formulas depend on $f$ and $f'''$, which we do not know. What matters is how they depend on $\epsmach$. To see
it, suppose that $f$ varies on a length scale $w$, in the units of $x$. By this we mean that each derivative is about
the previous one divided by $w$: $\abs{f'} \approx \abs{f} / w$ and $\abs{f'''} \approx \abs{f} / w^3$.\mnote{The
\emph{scale} $w$ of $f$: the distance in $x$ over which $f$ changes appreciably.} Putting these into the two formulas
gives
\begin{equation*}
  \hbest \approx 3^{1/3} \epsmach^{1/3} w,
  \qquad
  \frac{E(\hbest)}{\abs{f'}} \approx \epsmach^{2/3}.
\end{equation*}
In the second one, the constant $\frac{3}{2} \cdot 3^{-1/3}$ is close to 1 and has been dropped. For a \code{double},
$\epsmach^{1/3} = \num{\valRootEpsThreeDouble}$ and $\epsmach^{2/3} = \num{\valModelCentralTwoFirstDouble}$. So at
its best step the central difference keeps about \num{\valModelDigitsCentralTwoFirstDouble} of the 16 digits (Nocedal
and Wright,\cite{nocedal2006} Section 8.1). The same argument works for the forward difference, whose error is
$h \abs{f''} / 2$. It gives a best step near $\epsmach^{1/2} w$, and a relative error near
$\epsmach^{1/2} = \num{\valModelForwardOneFirstDouble}$: about 8 digits. That difference in accuracy, not the cost,
is why the central difference is the default.

The relative figure $\epsmach^{2/3}$ rests on an assumption: $\abs{f} \approx w \abs{f'}$. Look again at the rounding
term of \cref{eq:error-model}. It grows with $\abs{f}$, not with $\abs{f'}$. So a function whose values are large
compared with $w \abs{f'}$ loses accuracy in proportion. It loses about $\log_{10} (\abs{f} / (w \abs{f'}))$ digits.
For the nitrogen example, with $w = T_0$, the ratio is $\abs{\enth} / (T_0 \abs{\cp}) = \num{\valCpRelativeFactor}$,
close to 1. Suppose instead that the enthalpy had been referred to \SI{0}{\kelvin}, or had included an enthalpy of
formation, a constant that shifts every value. Then $\abs{\enth}$ would be larger, and the derivative proportionally
less accurate, although $\cp$ would be the same. A formation enthalpy of a few hundred kilojoules per mole, against
$T_0 \cp$ of a few tens, makes the ratio about ten. That costs about one digit. Near a point where $f'$ vanishes, the
relative error, which divides by $f'$, has no meaning at all. The absolute error, however, stays near
$\epsmach^{2/3} \abs{f} / w$.

On the nitrogen example, the measured curve reaches its lowest point, \num{\valCpBestRelError}, at
\SI{\valCpBestStep}{\kelvin}. That is below the model, because the rounding errors happen to cancel at that step. A
neighbouring step does several times worse. So what you should expect is the model's level, not the luckiest point.

\section{The default step}\label{sec:rounding-default}

The best step needs $f$ and its scale. The library does not know $f$, so it can compute neither $\hbest$ nor the
scale $w$. Instead it takes the formula of \cref{sec:rounding-best} without its constant, and it guesses the scale:
\begin{equation}
  h = \epsmach^{1/3} \scale, \qquad \scale = \max(\abs{x}, 1).
  \label{eq:default-step}
\end{equation}
The scale $\scale$ is the library's guess for $w$. For large $\abs{x}$, it makes the step relative: a fixed fraction
of $x$. The reason is that functions of a physical quantity often vary on a scale comparable to the size of that
quantity. The enthalpy, for example, varies with $T/\SI{1000}{\kelvin}$. Near zero, $\abs{x}$ says nothing about the
scale. There the step becomes absolute: $\epsmach^{1/3}$, in the units of $x$. This choice of scale is standard in
numerical optimisation (Dennis and Schnabel,\cite{dennis1996} Section 5.4).

For the nitrogen example, $\scale = \SI{1000}{\kelvin}$, and the default step is \SI{\valCpDefaultStep}{\kelvin}. The
relative error is \num{\valCpDefaultRelError}, which lies on the model line near the bottom of the V
(\cref{fig:v-curve}). Every stencil of the library has a default step of the same form,
$\epsmach^{1/(\aorder + \dorder)} \scale$ (\cref{ch:stencils,ch:orders}). For the central difference, $\aorder = 2$
and $\dorder = 1$, which gives back $\epsmach^{1/3} \scale$.

Dropping the constant $3^{1/3} \approx \num{\valStepConstantCentralTwoFirst}$ changes the step by that factor. When
the guess is right, $w = \scale$, the default step is smaller than the model's best step by that factor. This costs
little, and
the guess itself is the weaker link (\cref{sb:step-constant}).

\begin{sidebar}[label=sb:step-constant]{The error model behind the default step, and why its constant is dropped}
Take a rule of accuracy order $\aorder$ for the $\dorder$-th derivative. Its error model is
$E(h) = \errc h^{\aorder} \abs{f^{(\aorder + \dorder)}} + \ampl \epsmach \abs{f} / h^{\dorder}$. Here $\errc$ is the
error constant of the rule, and $\ampl$ the sum of its absolute weights (\cref{ch:stencils}). For the central
difference, $\errc = 1/6$, $\ampl = 1$ and $\dorder = 1$, and the model is \cref{eq:error-model}. The model is
smallest at
\begin{equation*}
  \hbest = \left( \frac{\dorder \ampl \epsmach \abs{f}}{\aorder \errc \abs{f^{(\aorder + \dorder)}}}
  \right)^{1/(\aorder + \dorder)},
\end{equation*}
and if $f$ varies on the scale $w$, this is $k \epsmach^{1/(\aorder + \dorder)} w$, with the constant
$k = (\dorder \ampl / (\aorder \errc))^{1/(\aorder + \dorder)}$. The default step leaves $k$ out, and puts the guess
$\scale$ in place of $w$. So when the guess is right, the default step equals $\hbest / k$. With $h = \hbest / k$,
the model error grows by the factor
\begin{equation*}
  \frac{E(\hbest / k)}{E(\hbest)} = \frac{\dorder k^{-\aorder} + \aorder k^{\dorder}}{\aorder + \dorder}.
\end{equation*}
For the central difference, $k = 3^{1/3} \approx \num{\valStepConstantCentralTwoFirst}$. Putting $\aorder = 2$ and
$\dorder = 1$ into the formula, the factor is $7 \cdot 3^{1/3} / 9 \approx \num{\valStepCostCentralTwoFirst}$. So
the default step loses about a tenth of the best accuracy. For the other first-derivative rules, $k$ lies between
that value and 2, the value for \code{Forward<1>}. The cost reaches \num{\valStepCostCentralTenFirst} for
\code{Central<10>}. For higher derivatives $k$ is larger: \num{\valStepConstantCentralTwoSecond} for
\code{Central<2>} and $f''$, and \num{\valStepConstantCentralTwoFourth} for $f^{(4)}$. So is the cost
(\cref{ch:orders}, \cref{sec:derivations-model}).

Why leave $k$ out at all? Because $k$ only matters when the guess $w = \scale$ holds, and the guess is the weaker
link. A step that is off by a factor 2 either way costs up to a factor \num{\valMisjudgeCostCentralTwoFirst} for the
central difference. That is more than leaving out $k$ costs. And a function that varies faster than $\scale$ wants a
smaller step anyway. The library's design record (\code{docs/deriv/PHASE4_API.md}) has nine test functions, which
vary faster than $\max(\abs{x}, 1)$. On them, half the default step did better for every order, $f'$ included. So the
library keeps one convention, $\epsmach^{1/(\aorder + \dorder)} \scale$, for every rule and every order, and leaves
the scale to the caller (\cref{sec:rounding-scale}). \Cref{ch:derivations} gives the derivation in full.
\end{sidebar}

There is one more limit on a small step: it must be large enough to move $x$ at all. The doubles near
\SI{1000}{\kelvin} are \SI{\valUlpAtOneThousand}{\kelvin} apart. So $T_0 + \SI{e-14}{\kelvin}$ rounds back to $T_0$,
and the two points of the central difference would coincide. The library checks for coinciding points. Rather than
divide by zero, it returns \code{error::non_finite}:
\needspace{7\baselineskip}
\snippet{03-rounding-step}{too-small}
The default step is far from this limit, because $\epsmach^{1/3} \scale$ is many spacings of $x$. A step that you
pass yourself may not be. The sidebar below (\cref{sb:ulp}) explains how finely $x$ is resolved, and how $x + h$ is
rounded.

\begin{sidebar}[label=sb:ulp]{Machine epsilon, ulps and the rounding of $x + h$}
The spacing of the numbers at $x$ is called its \emph{unit in the last place}, or ulp. It is
$\ulp(x) = \epsmach \cdot 2^{e}$ for $2^{e} \le \abs{x} < 2^{e+1}$. So $\ulp(1) = \epsmach$; just below 1 the spacing
is half of that; and $\ulp(1000) = \num{\valUlpAtOneThousand}$. Rounding to the nearest number errs by at most
$\ulp(x) / 2$. Relative to $x$, that is at most $\epsmach / 2$, the \emph{unit roundoff}. This book and the library's
rounding model use $\epsmach$ instead, which gives a factor 2 of safety (\cref{sec:fp-epsilon}).

Since $x + h$ is rounded, the point that the function sees is $\fl(x + h)$, not $x + h$. If $h$ is below half a
spacing of $x$, then $\fl(x + h) = x$, and the step has no effect. When this happens on both sides, both points round
back to $x$. No difference can be formed, and the library reports \code{error::non_finite}. Near a power of two,
however, the spacing differs on the two sides. Below $2^e$ the numbers are twice as dense as above it. So a step can
move $x$ downwards but not upwards. The central difference then divides by the half-spacing that remains. It returns
a finite value that is mostly rounding, without failing. The rules with more than two nodes check both $x + h$ and
$x - h$, and fail instead (\cref{sec:api-derivative}). Hence the documented requirement: a step well above
$\ulp(x)$.
\end{sidebar}

\section{When the default step is wrong: the scale of $f$}\label{sec:rounding-scale}

The default step assumes that the guess is right, $w = \scale$. In words, it assumes that $f$ varies on the scale
$\max(\abs{x}, 1)$, in the units of $x$. This section shows two common situations that break that assumption. The
remedy is the same for both.

\paragraph{A narrow feature at a large $\abs{x}$.} Take a function that changes from $-1$ to $1$ over a width of
about $w = 0.5$ around $x_0 = 1000$, such as $\tanh((x - x_0) / w)$. We differentiate it at $x = x_0 + 0.2$. You meet
this situation with a smoothed switch, or with a sharp transition in a variable far larger than its width. Here
$\scale = \num{1000.2}$, so the default step is \num{\valScaleDefaultStep}. That is small compared with the width, but
not small enough. The truncation error grows as $(h/w)^2$, not as $(h/\scale)^2$, and the relative error is
\num{\valScaleDefaultRelError}. \Cref{fig:step-scale} shows why: the default step sits far to the right of the bottom
of the V. The remedy is to put the true scale $w$ in place of the guess $\scale$:
\begin{equation*}
  h = \epsmach^{1/3} w,
\end{equation*}
which gives the step \num{\valScaleFitStep} and the error \num{\valScaleFitRelError}, at the bottom of the V. Ridders'
method with an initial step on the same scale, $w/4$, is just as good. It returns an error of
\num{\valScaleRiddersUserRelError}, with an estimate that covers it (\cref{ch:ridders}). In code, $\epsmach^{1/3}$ is
\code{root_epsilon<3, double>()}:\mnote{\code{nxx::math::root_epsilon<N, T>()} returns $\epsmach^{1/N}$ for the type
\code{T}, correct to within an ulp (\cref{sec:api-real}).}
\needspace{15\baselineskip}% language pass: keeps the listing whole, not split after its first line
\snippet{03-rounding-scale}{narrow-feature}

\begin{figure}
\sidecaption{The V for a narrow feature at a large $\abs{x}$: $\tanh((x - x_0)/w)$ with $x_0 = 1000$ and $w = 0.5$,
at $x = x_0 + 0.2$. The default step $\epsmach^{1/3} \scale$ sits far to the right of the minimum. The step
$\epsmach^{1/3} w$, built from the width of the feature, lands at the minimum. The default step guesses that $f$
varies on the scale $\scale = \max(\abs{x}, 1)$. Put the scale $w$ on which $f$ does vary in place of $\scale$.%
\label{fig:step-scale}}
\centering
% F4 fig:step-scale (Chapter 3, text width): the V of the central difference (Central<2>) for a narrow feature at a
% large |x|, tanh((x - x0) / w) with x0 = 1000 and w = 0.5 at x = x0 + 0.2 (DS4, data/step-scale.csv). The default
% step eps^(1/3) s lies far to the right of the minimum; the step eps^(1/3) w from the feature's width lands at it.
% Both steps and their errors come from \valScale{Default,Fit}{Step,RelError}.
\begin{tikzpicture}
\begin{axis}[eqtufte, width=95mm, height=45mm, xmode=log, ymode=log,
    xlabel={step $h$}, ylabel={relative error},
    xmin=1e-12, xmax=10, ymin=1e-12, ymax=1,
    xtickten={-12,-10,-8,-6,-4,-2,0}, ytickten={-12,-9,-6,-3,0}]
  \addplot[black!75, mark=*] table[col sep=comma, x=h, y=error, comment chars={\#}] {data/step-scale.csv};
  % the default step eps^(1/3) max(|x|, 1)
  \draw[muted, thin] ({axis cs:\valScaleDefaultStep,1}|-{rel axis cs:0,0})
    -- ({axis cs:\valScaleDefaultStep,1}|-{rel axis cs:0,1})
    node[above, font=\sffamily\scriptsize, inner sep=1.5pt] {default step $\varepsilon^{1/3} s$};
  \fill[muted] (axis cs:\valScaleDefaultStep,\valScaleDefaultRelError) circle (1.4pt);
  % the step from the scale w on which f varies
  \draw[accent, thin] ({axis cs:\valScaleFitStep,1}|-{rel axis cs:0,0})
    -- ({axis cs:\valScaleFitStep,1}|-{rel axis cs:0,1})
    node[above, font=\sffamily\scriptsize, inner sep=1.5pt] {$\varepsilon^{1/3} w$};
  \fill[accent] (axis cs:\valScaleFitStep,\valScaleFitRelError) circle (1.4pt);
\end{axis}
\end{tikzpicture}

\end{figure}

\paragraph{A small $\abs{x}$ in your units.} Think of a mole fraction of $10^{-6}$, a strain of $10^{-4}$, or a time
constant in seconds of a fast process. All of them have $\scale = 1$, so the default step is $\epsmach^{1/3}$. That
is a sizeable fraction of the variable, or larger than it. Take $y \ln y$, which appears in every entropy of mixing, at
$y = 10^{-6}$. The default step is \num{\valSmallDefaultStep}, which exceeds $y$. So the central difference samples
$y - h < 0$, where the logarithm is not defined, and the call fails with \code{error::non_finite}. With the step
$\epsmach^{1/3} y = \num{\valSmallFitStep}$, the relative error is \num{\valSmallFitRelError}. That is as accurate as
at any ordinary point (\cref{fig:step-scale-small}). The lesson: the 1 in $\max(\abs{x}, 1)$ is in your units. It
means 1 in whatever unit you measure $x$, so your choice of unit changes the default step.
\snippet{03-rounding-scale}{small-x}
Here the failure is loud, because the logarithm returns NaN below zero. Now think of a function that is defined
across the bound but meaningless there, such as a fit sampled outside its range. It would return a wrong value
without complaint. The one-sided rules of \cref{ch:stencils} are the tool for such an edge.

Units matter at large $\abs{x}$ too. A property near room temperature varies on the scale of the absolute
temperature. But at \SI{20}{\degreeCelsius} the guess is $\scale = 20$ in degrees Celsius, against 293 in kelvin. The
default step is then about 15 times too small. On the rounding branch, where the error grows as $1/h$, that costs
about one digit.

\begin{marginfigure}
\centering
% F4b fig:step-scale-small (Chapter 3, margin): the V of the central difference (Central<2>) for y ln y at
% y = 1e-6 (DS4b, data/step-scale-small.csv). Steps above y sample below zero, where the function is not finite, so
% the curve stops there (nan in the CSV) and the default step eps^(1/3) gets no value; the step eps^(1/3) y lands
% near the minimum. Steps, status and error from \valSmall{DefaultStep,DefaultStatus,FitStep,FitRelError}.
\begin{tikzpicture}
\begin{axis}[eqtufte, width=36mm, height=32mm, xmode=log, ymode=log,
    xlabel={step $h$}, ylabel={relative error},
    xmin=1e-20, xmax=1e-5, ymin=1e-12, ymax=1e-2,
    xtickten={-20,-15,-10,-5}, ytickten={-12,-8,-4},
    unbounded coords=jump]
  \addplot[black!75, mark=*] table[col sep=comma, x=h, y=error, comment chars={\#}] {data/step-scale-small.csv};
  \draw[muted, thin] ({axis cs:\valSmallDefaultStep,1}|-{rel axis cs:0,0})
    -- ({axis cs:\valSmallDefaultStep,1}|-{rel axis cs:0,1})
    node[anchor=south east, font=\sffamily\scriptsize, inner sep=1pt, align=right] {default step:\\ \valSmallDefaultStatus};
  \draw[accent, thin] ({axis cs:\valSmallFitStep,1}|-{rel axis cs:0,0})
    -- ({axis cs:\valSmallFitStep,1}|-{rel axis cs:0,1});
  \node[accent, font=\sffamily\scriptsize, anchor=east, inner sep=1.5pt]
    at ({axis cs:\valSmallFitStep,1}|-{rel axis cs:0,0.72}) {$\varepsilon^{1/3} y$};
  \fill[accent] (axis cs:\valSmallFitStep,\valSmallFitRelError) circle (1.2pt);
\end{axis}
\end{tikzpicture}

\caption{The V of \code{Central<2>} for $y \ln y$ at $y = 10^{-6}$. Steps above $y$ cross zero, so the curve stops
there and the default step fails. The step $\epsmach^{1/3} y$ gives the expected accuracy. The 1 in
$\max(\abs{x}, 1)$ is in your units.\label{fig:step-scale-small}}
\end{marginfigure}

\paragraph{Finding the scale.} The scale $w$ can usually be read off the model that $f$ computes. For a transition,
it is the width of the transition. For a singularity or a bound at $a$, such as in $1/(x - a)$ or $\ln(x - a)$, it is
the distance $\abs{x - a}$. For an Arrhenius factor $\exp(-E_a / (\Rgas T))$, the temperature dependence of a rate
constant, it is $\Rgas T^2 / E_a$, the change in $T$ over which the exponent changes by about 1. That is much smaller
than $T$ for a large activation energy. For a smooth interpolant of a table, such as a cubic spline, it is the node
spacing, although a spline is better differentiated by its own derivative (\cref{sec:using-not}).

A rough value is enough. Compared with the model's best step, a step that is off by a factor of two either way costs
\code{Central<2>} at most a factor \num{\valMisjudgeCostCentralTwoFirst} in error. If the scale is hard to judge, you
can use Ridders' method with an initial step a few times smaller than the scale. It chooses its own step, and reports
an error estimate (\cref{ch:ridders}).

\section{When $f$ is less accurate than $\epsmach$}\label{sec:rounding-noise}

So far the error model has assumed that each value of $f$ is accurate to about $\epsmach$. Many engineering functions
are not. Think of the output of an iterative solver stopped at a relative tolerance of $10^{-8}$; of a flash
calculation, which finds the phases of a mixture in equilibrium; or of a simulation that iterates to convergence.
Such an output is accurate only to the tolerance. Its error also changes irregularly from one argument to the next.
An error of this kind is called \emph{noise}. We call its relative size the \emph{noise level} $\noise$. It takes the
place of $\epsmach$ in the rounding term of \cref{eq:error-model},
\begin{equation*}
  E(h) = \frac{h^2}{6} \abs{f'''} + \frac{\noise \abs{f}}{h},
\end{equation*}
and the argument of \cref{sec:rounding-best} carries over, with $\noise$ for $\epsmach$. The best step is about
$\noise^{1/3} w$, and the best relative accuracy about $\noise^{2/3}$.\cite{more2012} So a solver converged to
$10^{-8}$ gives about five correct digits, not \num{\valHeadlineFirstDigits}. And it gives them only at a step
$(\noise / \epsmach)^{1/3}$ times the default.

Two cautions on $\noise$. First, a solver's tolerance is the noise level only when its convergence test is on the
output itself, relative to the output. Some solvers test something else. A flash may converge a residual in its own
units, such as an energy balance or the equality of fugacities; a residual is the amount by which an equation fails
to hold. A simulator may converge its tear streams, the recycle streams whose values it guesses and then iterates.
For such solvers the tolerance is not $\noise$. Measure $\noise$ instead, as shown below. Second, $f$ may itself be
computed by cancellation, as a difference of large terms. Then its rounding error is $\epsmach$ times the largest
intermediate value, that is $k \epsmach \abs{f}$ with $k \gg 1$. Treat $k \epsmach$ as the noise level.

Note that a smooth but inaccurate $f$ is not noisy. The Shomate fit for nitrogen describes nitrogen to only about
four digits. Yet its derivative, which is the derivative of the fit, is accurate to $\epsmach^{2/3}$.

\begin{marginfigure}
\centering
% F17 fig:noise (Chapter 3, margin): the V of the central difference (Central<2>) on exp(x) at x = 1, clean and with
% the values perturbed by a relative noise eta = 1e-10 and 1e-8 (DS19, data/step-error-noise.csv, the perturbation
% of code/noise.hpp). The rounding branch lifts to eta |f| / h and the minimum moves right and up. The default step
% (\valRootEpsThreeDouble, since s = 1) and the step eta^(1/3) for eta = 1e-8 with its error (\valNoiseEightFit...)
% are marked. Steps 1 to 1e-8 (on a log axis the domain restriction is in natural logarithms: ln(1e-8) = -18.42).
\begin{tikzpicture}
\begin{axis}[eqtufte, width=36mm, height=32mm, xmode=log, ymode=log,
    xlabel={step $h$}, ylabel={relative error},
    xmin=1e-8, xmax=1, ymin=1e-12, ymax=1,
    xtickten={-8,-4,0}, ytickten={-12,-8,-4,0},
    restrict x to domain=-18.43:0.01]
  \addplot[muted, mark=*] table[col sep=comma, x=h, y=clean, comment chars={\#}] {data/step-error-noise.csv};
  \addplot[black!75, mark=*] table[col sep=comma, x=h, y=noise_ten, comment chars={\#}] {data/step-error-noise.csv};
  \addplot[accent, mark=*] table[col sep=comma, x=h, y=noise_eight, comment chars={\#}] {data/step-error-noise.csv};
  \draw[muted, thin] ({axis cs:\valRootEpsThreeDouble,1}|-{rel axis cs:0,0})
    -- ({axis cs:\valRootEpsThreeDouble,1}|-{rel axis cs:0,1})
    node[above, font=\sffamily\scriptsize, inner sep=1.5pt] {default};
  \fill[accent] (axis cs:\valNoiseEightFitStep,\valNoiseEightFitRelError) circle (1.2pt);
  \node[accent, font=\sffamily\scriptsize, anchor=south west, inner sep=1pt] at (axis cs:2e-5, 1e-3)
    {$\eta = 10^{-8}$};
  \node[black!75, font=\sffamily\scriptsize, anchor=north west, inner sep=1pt] at (axis cs:1e-7, 1e-6)
    {$10^{-10}$};
  \node[muted, font=\sffamily\scriptsize, anchor=south west, inner sep=1pt] at (axis cs:1e-7, 1e-9)
    {clean};
\end{axis}
\end{tikzpicture}

\caption{The V of \code{Central<2>} on $\mathrm{e}^x$ at 1, with the values perturbed by a relative noise
$\noise = 10^{-10}$ and $10^{-8}$, and the clean curve for comparison. The rounding branch lifts to
$\noise \abs{f} / h$, and the minimum moves right and up. The dot marks the step $\noise^{1/3}$ for
$\noise = 10^{-8}$.\label{fig:noise}}
\end{marginfigure}

\Cref{fig:noise} measures this on $\mathrm{e}^x$ at $x = 1$. We multiply its values by $1 + \noise u$, with $u$
drawn uniformly from $[-1, 1]$. The same $x$ always gets the same $u$, as it would from a converged solver. For
$\noise = 10^{-8}$, the default step returns an error of \num{\valNoiseEightDefaultRelError}. The step
$\noise^{1/3} = \num{\valNoiseEightFitStep}$ gives \num{\valNoiseEightFitRelError}, close to the model's
$\noise^{2/3} = \num{\valNoiseEightModelBestRelError}$ (\code{code/03-rounding-noise.cpp}). So when $f$ is noisy,
pass the step yourself: $h = \noise^{1/3} \scale$, with $\noise$ the solver's tolerance or your best estimate of the
accuracy of $f$. In code, this is \code{derivative(f, x, std::cbrt(eta) * std::max(std::abs(x), 1.0))}.

If you do not know the noise level, you can measure it. Evaluate $f$ at a handful of points, spaced far more closely
than the scale on which $f$ varies. Then form their successive differences, as in a difference table: the
differences of neighbouring values, then the differences of those, and so on. The $k$-th differences of the smooth
part of $f$ shrink like $h^k$. So at such a spacing, the higher differences show the noise alone, and their spread
estimates $\noise \abs{f}$ up to a known factor. That factor is the growth of the differences of independent noise:
their $k$-th differences are $\sqrt{(2k)! / k!^2}$ times larger than the noise itself, and the recipe divides this
factor out. The recipe is a simplified form of Mor\'e and Wild's procedure,\cite{more2011} and the book's example
header keeps it:
\snippetheader{noise}{noise-estimate}
We apply it to the friction factor of \cref{sec:using-solver}, which is computed by bisection to a relative tolerance
of $10^{-10}$. We use eight points, spaced by the default step:
\needspace{13\baselineskip}% language pass: keeps the listing whole, not split after its first line
\snippet{03-rounding-noise-estimate}{noise-estimate-use}
The recipe reports a noise level of \num{\valNoiseEstimateFriction} relative to $f$. It settles at order
\num{\valNoiseEstimateFrictionK}, and uses \num{\valNoiseEstimateEvals} evaluations of $f$. The actual root-mean-square
deviation from the exact root is \num{\valNoiseEstimateFrictionActual}. On the noisy $\mathrm{e}^x$ above, it reports
\num{\valNoiseEstimateNoisy}. That is near $\noise / \sqrt{3}$, the root-mean-square size of $\noise u$, for
$\noise = \num{\valNoiseEstimateNoisyEta}$. On the clean $\mathrm{e}^x$, it reports \num{\valNoiseEstimateClean}.

If the estimates of no three successive orders agree, the recipe reports $k = 0$. This means that the spacing is too
large: within eight points, the smooth part does not drop below the noise. Then retry with a tenth of the spacing, or
fall back to the solver's tolerance, as the listing does. A spacing on the scale of $f$ itself can also make the
orders agree by accident, at a wrong level. (The full procedure tests signs and adjusts the spacing.) So use the
default step, check $k$, and distrust a level near $\epsmach$ for a solver. The header,
\code{docs/deriv-book/code/noise.hpp}, is under the repository's MIT licence.

Ridders' method needs care with noise too. Its error estimate assumes that $f$ is accurate to about $\epsmach$. So
for a noisy $f$, the estimate can be optimistic, that is, smaller than the actual error. Over
\num{\valNoiseRiddersCases} points in $[0.2, 3]$, the estimate covered the actual error in
\SI{\valNoiseCleanRiddersCoverage}{\percent} of the cases for the clean $\mathrm{e}^x$. For $\noise = 10^{-8}$ it
covered it in only \SI{\valNoiseEightRiddersCoverage}{\percent}, and the worst error was
\num{\valNoiseEightRiddersWorstRatio} times its estimate. \Cref{sec:using-solver} differentiates the friction factor
with the noise-adjusted step of this section.

\section{Float, double and more digits}\label{sec:rounding-types}

Everything in this chapter holds for any number type, each with its own $\epsmach$. A \code{float} has $\epsmach =
\num{\valEpsFloat}$, so it reaches only about $\epsmach^{2/3} = \num{\valModelCentralTwoFirstFloat}$. On
$\mathrm{e}^x$ at 1, its default step gives \num{\valFloatDefaultRelError} (\cref{fig:float-double}). A 50-digit type,
with $\epsmach = \num{\valEpsMultiprecision}$, reaches far beyond \code{double} (\cref{sec:using-types}).

\begin{marginfigure}
\centering
% F5 fig:float-double (Chapter 3, margin): the V of the central difference (Central<2>) on exp(x) at x = 1 in float
% and in double (DS3, data/step-error-float.csv). Each floor sits at about eps^(2/3) of its type, drawn as a dashed
% guide from \valModelCentralTwoFirst{Float,Double}. Steps 1 to 1e-8 (on a log axis the domain restriction is in
% natural logarithms: ln(1e-8) = -18.42).
\begin{tikzpicture}
\begin{axis}[eqtufte, width=36mm, height=32mm, xmode=log, ymode=log,
    xlabel={step $h$}, ylabel={relative error},
    xmin=1e-8, xmax=1, ymin=1e-12, ymax=1,
    xtickten={-8,-4,0}, ytickten={-12,-8,-4,0},
    unbounded coords=jump, restrict x to domain=-18.43:0.01]
  \draw[accent, densely dashed, line width=0.5pt] ({rel axis cs:0,0}|-{axis cs:1,\valModelCentralTwoFirstFloat})
    -- ({rel axis cs:1,0}|-{axis cs:1,\valModelCentralTwoFirstFloat})
    node[below left, font=\sffamily\scriptsize, inner sep=1pt] {$\varepsilon_{\mathrm{float}}^{2/3}$};
  \draw[accent, densely dashed, line width=0.5pt] ({rel axis cs:0,0}|-{axis cs:1,\valModelCentralTwoFirstDouble})
    -- ({rel axis cs:1,0}|-{axis cs:1,\valModelCentralTwoFirstDouble})
    node[above left, font=\sffamily\scriptsize, inner sep=1pt] {$\varepsilon_{\mathrm{double}}^{2/3}$};
  \addplot[muted, mark=*] table[col sep=comma, x=h, y=error_float, comment chars={\#}] {data/step-error-float.csv};
  \addplot[black!75, mark=*] table[col sep=comma, x=h, y=error_double, comment chars={\#}] {data/step-error-float.csv};
  \node[muted, font=\sffamily\scriptsize, anchor=south west] at (axis cs:5e-7, 0.12) {float};
  \node[black!75, font=\sffamily\scriptsize, anchor=south west] at (axis cs:1.2e-8, 1e-8) {double};
\end{axis}
\end{tikzpicture}

\caption{The V of \code{Central<2>} for \code{float} and \code{double}, on $\mathrm{e}^x$ at 1. The floor of each
curve lies at or below $\epsmach^{2/3}$ of its type (dashed). The rounding scatter reaches a decade or more below
that level.\label{fig:float-double}}
\end{marginfigure}

\section{A checklist for the step}\label{sec:rounding-checklist}

This section gathers what the chapter has shown into six questions. Answer them before you accept the default step.
\begin{enumerate}
  \item \emph{On what scale does $f$ vary, in the units of $x$?} If the scale is not about $\max(\abs{x}, 1)$, pass
    $h = \epsmach^{1/3} w$ for \code{Central<2>}, where $w$ is that scale (\cref{sec:rounding-scale}).
  \item \emph{Is there a bound within a step of $x$}, such as zero for a mole fraction or the end of a fit's range?
    If so, use a smaller step, or a one-sided rule (\cref{ch:stencils}).
  \item \emph{How accurate is $f$?} Suppose that $f$ is a solver output, or otherwise noisy, at a level $\noise$
    above $\epsmach$. Then pass $h = \noise^{1/3} \scale$, and expect a relative error of about $\noise^{2/3}$. If
    you do not know $\noise$, measure it (\cref{sec:rounding-noise}).
  \item \emph{Is $\abs{f}$ much larger than $\scale \abs{f'}$}, as for an enthalpy with a formation term or an
    absolute pressure with a small slope? If so, expect to lose about $\log_{10} (\abs{f} / (\scale \abs{f'}))$
    digits. If you can, subtract a constant reference from $f$ before you differentiate; a constant does not change
    $f'$ (\cref{sec:rounding-best}).
  \item \emph{Which method and order?} A step chosen for one rule is wrong for another (\cref{ch:orders}).
  \item \emph{Unsure?} Use Ridders' method, which chooses its own step and estimates its error (\cref{ch:ridders}).
    Or repeat a stencil at $2h$ and $h/2$ (\cref{sec:using-check}).
\end{enumerate}
